As intelligence systems start playing important roles in our daily life, 
ethics in artificial intelligence research has attracted significant interest.
It is known that big-data technologies sometimes inadvertently worsen 
discrimination due to implicit biases in data~\cite{PPMHZ14}. 
Such issues have been demonstrated in various learning systems, including
online advertisement systems~\cite{sweeney2013discrimination}, word embedding
models~\cite{bolukbasi2016man,CBN17}, online news~\cite{ross2011women}, 
web search~\cite{kay2015unequal}, and credit score~\cite{hardt2016equality}.
Data collection biases have been discussed in the context of creating image corpus~\cite{misra2016seeing,van2016stereotyping} and text corpus~\cite{gordon2013reporting,van2010extracting}. In contrast, we show that given a gender biased corpus, structured models such as conditional random fields, amplify the bias. 

The effect of the data imbalance can be easily detected and fixed when the prediction task is simple. 
For example, when classifying binary data with unbalanced labels (i.e., samples in the majority class dominate the dataset), a classifier trained exclusively to optimize accuracy learns to always predict the majority label, as the cost of making mistakes on samples in the minority class can be neglected. 
Various approaches have been proposed to make a ``fair'' binary classification~\cite{barocas2014big,dwork2012fairness,feldman2015certifying,zliobaite2015survey}.
For structured prediction tasks the effect is harder to quantify and we are the first to propose methods to reduce bias amplification in this context. 

Lagrangian relaxation and dual decomposition techniques have been widely used in NLP tasks~(e.g., \cite{sontag2011introduction,Rush2012,chang2011exact,peng2015dual}) for dealing with instance-level constraints. Similar techniques~\cite{ChangSuKe13,dalvi2015constrained} have been applied in handling  corpus-level constraints for semi-supervised  multilabel classification. 
In contrast to previous works aiming for improving accuracy performance, we incorporate corpus-level constraints for reducing gender bias.  